\documentclass{article}
\usepackage{setspace}
\begin{document}
%Project Summary.

%The Project Summary is a one page, single-spaced summary of the proposed
%activity. The Project Summary should be a self-contained
%description of the proposed research. It should include a statement of
%objectives and methods to be employed. It must address the scientific and/or
%technical merit of the project, for example, the influence that the results
%might have on the direction, progress and thinking in relevant scientific or
%engineering fields.

%It must also address the appropriateness of the proposed methods, and the
%logic and feasibility of the research approach.  The Project Summary should be
%informative to persons working in the same or related fields and, insofar as
%possible, understandable to a scientifically or technically literate lay
%reader.
%%%%%%%%%%%%%%%%%%%%%%%%%%%%%%%%%%%%%%%%%%%%%%%%%%%%%%%%%%%%%%%%%%%%%%%%%%%%%%%%

\begin{singlespace}
  

%The automation of tasks that are dangerous, inconvenient, or otherwise undesirable for humans are one of the main benefits of intelligent and programmable robots. 
Environmental monitoring is a field that could benefit greatly from efforts in automation as the dynamic and spatially transient properties of natural phenomena can be difficult or costly to carry out via traditional human intensive methods. 

We propose research into the design of small environmentally adaptable robotic platforms to expand available methods for environmental monitoring with the objective of increasing the quality and quantity of data collection through greater automation. 

Towards that general goal we identify certain traits to focus on. Multi-domain function and navigation, to expand the scope of measurements. Field deployability, platforms should be mechanically, chemically, and electrically robust to function in real-world scenarios. Moreover, in order to promote use and adoption these platforms should be an aid to existing scientists in the field, therefore, they should be easy for them to use and program in the field. 

We propose developing a land-water robot for assisting water quality measurements and oil slick detection to depths of up to 10 meters and with the capacity to carry sensors to its native proprioceptive sensors in backpack modules. One such sensor is a water sampling backpack. 

We also propose the development of an amphibious powertrain for small robotic platforms. This method of locomotion would allow a platform to perform on land and in water while being robust to field conditions, increasing the robot's adaptibility. Screwdrive locomotion is relatively unexplored method of locomotion, especially for robotics, however it is uniquely positioned to fulfill the amphibious task role while being simple in design. 

Prior research on large military tank-size screwdrive propulsion systems indicates that optimal screw parameters, such as blade height and pitch, differ for land and water regimes, therefore, characterization at this scale is necessary. Furthermore, we propose the development of a variable-pitch screw to be used in this system both for characterization purposes and for field deployment. One way to achieve this would be the use of origami or kirigami to induce buckling along predefined creases where actuation varies the pitch of the screw blade in a manner akin to the way an accordion expands and contracts. In order to match the variable screws to their terrain, a terrain classification system is also proposed.  

Finally, noting the benefits of land-water adaptability to the available mission set, we propose development of a land-water-air platform that incorporates principles from the land-water platform, the amphibious powertrain, and adds flight capabilities. This new platform, to the author's knowledge, would be the first of its kind. The addition of flight would allow quicker and easier deployment and retrieval while allowing sensing in spaces not previously available while maintaining land-water capabilities.

%Finally, we propose the development of a land-water-air platform that incorporates principles from the Aquapod, the screwdrive, and adds flight. This new platform, to the author's knowledge, would be the first of its kind. This platform would allow quicker deployment using flight to overcome obstacles while still being able to perform measurement tasks. 

\end{singlespace}

\end{document}